% Afgeleiden. De zeilwagen van Simon Stevin (1600): het snelste voertuig van zijn tijd, van Scheveningen
% naar Petten in twee uur. Snelheid op een ogenblik.
% Motief: koorden door P die samenvallen met de raaklijn, wier richting de afgeleide is.
\def\omslagbestand{omslag-zeilwagen.jpg}
\def\omslagbron{Anoniem naar Willem van Swanenburg en Jacques de Gheyn II, \textit{De zeilwagen van Simon Stevin}, ca. 1640-1649 (publiek domein)}
\def\omslagrand{10.4 + 6*(\x/21)^5}
\def\omslagmotief{%
  \randlijn
  \randwaarde{13}{\yp}%
  \pgfmathsetmacro{\rico}{30/21*(13/21)^4}%
  \foreach \q in {20.5,18.5,16}{%
    \randwaarde{\q}{\yq}%
    \pgfmathsetmacro{\yeind}{\yp+(\yq-\yp)/(\q-13)*(\q+.5-13)}%
    \draw[omslaglijn=.9pt,dashed] (13,\yp) -- (\q+.5,\yeind);
    \fill[omslagstip] (\q,\yq) circle (1.2mm);
  }
  \draw[fel,line width=1.8pt] (9,{\yp-4*\rico}) -- (21,{\yp+8*\rico});
  \fill[omslagstip] (13,\yp) circle (1.7mm);
  \node[omslaglabel,anchor=south east] at (12.85,{\yp+.12}) {$P$};
  \node[text=white,font=\small,inner sep=1pt,anchor=north east] at (20.6,{\yp+7.6*\rico-.2}) {raaklijn};
}
